\section{Introduction}

Text embedding models induce a geometry over natural language data. Similarity, clustering, semantic transfer, and retrieval behavior are downstream manifestations of this geometry: useful embeddings organize examples into neighborhoods, regions, and continuous similarity structures that reflect meaning. Strong recent embedding models, including LLM-based encoders, improve this organization but are costly to run at scale. This motivates compact student encoders that inherit the semantic structure of a stronger teacher without reproducing every teacher coordinate or internal state.

Knowledge distillation offers a practical route to this compression~\citep{hinton2015distilling}. In language representation learning, distillation has been applied through output matching, hidden-state alignment, self-attention transfer, and compact pretrained architectures~\citep{sanh2019distilbert,jiao2019tinybert,wang2020minilm,xu2024survey}. For embedding models, common objectives include pointwise alignment between student and teacher embeddings, contrastive learning with teacher-derived positives, and relational losses that match pairwise similarities within a mini-batch. These objectives are attractive because they are simple and efficient. They also fit naturally into stochastic training.

However, embedding models are evaluated through geometric properties that these losses only indirectly supervise. Class-level transfer depends on whether semantic regions remain separable, pairwise decisions depend on local boundaries, and similarity tasks depend on calibrated distances. A pointwise loss asks the student to reproduce teacher coordinates, even when a compact student may not have enough capacity to match the teacher's absolute geometry. A mini-batch relational loss improves on this by comparing local similarity patterns, but it sees only the examples that happen to co-occur in a batch. Many relations that matter for representation transfer are neither pointwise nor confined to a single batch: samples may be connected through multi-hop semantic paths, and useful regions often emerge from the accumulated geometry of the support set.

This mismatch motivates the central question of this work: can a compact student inherit the teacher-induced diffusion geometry of a semantic manifold, rather than only matching individual vectors or mini-batch relations? The intuition is simple. If the teacher's embedding space defines a semantic manifold over a support set, then random walks on a teacher-induced neighborhood graph reveal how semantic mass flows from each example to nearby examples, broader local regions, and higher-order communities.

Our contributions are:
\begin{itemize}
    \item We formulate text embedding distillation as preserving the teacher-induced semantic manifold and introduce Manifold Diffusion Distillation, which transfers multi-scale Markov diffusion profiles and low-frequency spectral structure into a compact student encoder.
    \item We provide finite-support analysis connecting the MDD objective to diffusion-geometry preservation, including multi-scale barycenter matching, candidate-truncation error, low-frequency topology preservation, and the coverage gap of batch-local relational objectives.
    \item We evaluate the method through geometry-sensitive frozen embedding probes covering class-region separability, pairwise semantic discrimination, and continuous similarity calibration.
\end{itemize}

MDD first caches teacher embeddings for an unlabeled support set and constructs a mutual $k$-nearest-neighbor graph using teacher cosine similarity. From this graph, it defines a Markov transition matrix and precomputes diffusion targets corresponding to one-step, two-step, and longer-range semantic random walks. During training, each anchor example is paired with a candidate set containing high-diffusion neighbors, teacher hard negatives, and random negatives. The student then learns a neighborhood distribution over the same candidates and minimizes a weighted KL divergence to the teacher diffusion targets.

MDD adds two stabilizing signals. First, a spectral manifold objective projects student embeddings into low-dimensional spectral coordinates computed from the teacher graph Laplacian, encouraging the student to preserve global semantic topology rather than only local neighbors. Second, a lightweight pointwise anchor aligns projected student embeddings with cached teacher embeddings, acting as a boundary condition rather than the main training signal. A simple curriculum introduces diffusion scales gradually: the student first learns local neighborhoods, then semantic regions, and finally longer-range diffusion with spectral anchoring.
